% EM structure and teaching revision 20260928
\subsection{测量加载与电功率的局部核对}
理想电压表不取电流，理想电流表不产生压降；实际仪表会改变被测状态。电阻 $R_x$ 与内阻 $R_V$ 的电压表并联时，若电流表测的是总电流，则
\[
 R_{\text{读}}=\frac UI=\frac{R_xR_V}{R_x+R_V}<R_x.
\]
相反，若电压表跨接 $R_x$ 与电流表内阻 $R_A$ 的串联组合，而电流表只测流过 $R_x$ 的电流，则 $R_{\text{读}}=R_x+R_A$。先画出仪表端点，再分析电流究竟流过哪些支路，误差方向便不需死记。
场描述给出局部功率密度 $p=\boldsymbol j\cdot\boldsymbol E$。对欧姆介质，$p=j^2/\sigma=\sigma E^2\ge0$；积分后与端口功率 $UI$ 相等。以长直均匀电阻为例，$j=I/S$、$E=U/\ell$，故 $\int p\dd V=UI$。
这里的非负性只适用于耗散电阻。电源支路可以输出功率，其端口 $UI$ 的符号取决于是否采用电流流入电压正端的关联参考方向。

